\section{Conclusion}\label{sec:conclusion}

We showed that context retrieval---not model inference---is the dominant
cost of context-aware maritime trajectory prediction, and that it is a
previously unrecognised spatial database workload. M-CTX recasts it as
such and accelerates it by two to three orders of magnitude with provable
correctness and bit-exact downstream equivalence. Its algorithmic core,
BR-LZ, is the first one-curve learned spatial index with a recall-completeness
proof and a segment-local extent that is provably tighter than the
expansion competitors require; its systems contributions---a $163\times$
linear-time SDF engine and a $64\times$ storage co-design---turn a
multi-CPU-day, multi-machine pipeline into a single-machine one. Across
four real maritime regions, nine baseline systems, scaling to $40$\,M
features, and $10^{7}$-record streaming, M-CTX preserves four pretrained
models' predictions byte-for-byte while cutting context construction from
hours to minutes.

The lesson generalises beyond maritime analytics: the retrieval workloads
hidden inside large spatiotemporal learning pipelines are first-class
database problems, and learned-index machinery---once equipped with the
recall guarantees those workloads demand---transfers cleanly to them. We
expect the BR-LZ recall theorem, the storage--accuracy co-design, and the
bit-exact drop-in discipline to carry directly to autonomous driving and
urban mobility, where context retrieval is the same un-indexed
bottleneck. M-CTX is released as an open drop-in replacement for the
standard context pipeline.
